\showpoints

\fullwidth{You may use a calculator to solve the following problem but you must \textbf{show your work} in the space provided.  You may \textbf{round to 3 decimal places after each step \emph{or} after the final calculation}. Any other rounding will result in a 3 point deduction.}

\newcommand{\genotype}[1]{$#1$}

\question[15]
Impatiens is a type of flower that has three flower colors. Red color is produced by \genotype{R_1R_1} individuals. Pink color is produced by \genotype{R_1R_2} individuals. White color is produced by \genotype{R_2R_2} individuals. A random sample of 248 flowers found that 42 were red, 175 were pink, and 31 were white. Calculate the two observed allele frequencies and the three observed genotype frequencies. \emph{Show whether the population is in Hardy-Weinberg equilibrium.}

The final values may sum to 0.999, 1.000 or 1.001 due to rounding error. (2 pts each plus 5 points for showing your work.)

\begin{multicols}{2}

\begin{tabular}{lC{1.1in}}
\toprule
&	Frequency \\
\midrule
 & \\[0.5em]
\genotype{R_1}		& \ifprintanswers \textbf{0.522} \else \rule{1in}{0.4pt} \fi\\[2em]
\genotype{R_2}			& \ifprintanswers \textbf{0.478} \else \rule{1in}{0.4pt} \fi\\[2em]
\genotype{R_1R_1}		& \ifprintanswers \textbf{0.169} \else \rule{1in}{0.4pt} \fi\\[2em]
\genotype{R_1R_2}		& \ifprintanswers \textbf{0.706} \else \rule{1in}{0.4pt} \fi\\[2em]
\genotype{R_2R_2}		& \ifprintanswers \textbf{0.125} \else \rule{1in}{0.4pt} \fi\\
\bottomrule
\end{tabular}

\columnbreak

Be sure to show whether the population is in Hardy-Weinberg equilibrium and \emph{show all work.}

\end{multicols}

\ifprintanswers \textbf{Sums to 1.0.} \fi 
